\documentclass[11pt]{article}
% LaTeX-only set-up (the web build does not read this file).
\usepackage[T1]{fontenc}
\usepackage[utf8]{inputenc}
\usepackage{lmodern}
\usepackage{amsmath,amssymb,amsthm}
\usepackage[a4paper,margin=1.1in]{geometry}
\usepackage{microtype}
\usepackage{booktabs}
\usepackage{enumitem}
\usepackage{newunicodechar}
\usepackage[hidelinks]{hyperref}
\setlist{itemsep=0.2em,topsep=0.4em}
% Statements carry their own explicit label through \head, so nothing is numbered automatically.
\newenvironment{statement}{\par\addvspace{\medskipamount}\noindent\ignorespaces}{\par\addvspace{\medskipamount}}
\newenvironment{frontnote}{\par\addvspace{\medskipamount}\small\noindent\ignorespaces}{\par\addvspace{\medskipamount}}
\renewcommand{\qedsymbol}{$\blacksquare$}
% Unicode that may appear inside verbatim and \texttt (Lean code)
\newunicodechar{α}{\ensuremath{\alpha}}\newunicodechar{β}{\ensuremath{\beta}}
\newunicodechar{γ}{\ensuremath{\gamma}}\newunicodechar{δ}{\ensuremath{\delta}}
\newunicodechar{μ}{\ensuremath{\mu}}\newunicodechar{ω}{\ensuremath{\omega}}
\newunicodechar{τ}{\ensuremath{\tau}}\newunicodechar{λ}{\ensuremath{\lambda}}
\newunicodechar{Π}{\ensuremath{\Pi}}\newunicodechar{Σ}{\ensuremath{\Sigma}}
\newunicodechar{∀}{\ensuremath{\forall}}\newunicodechar{∃}{\ensuremath{\exists}}
\newunicodechar{∧}{\ensuremath{\wedge}}\newunicodechar{∨}{\ensuremath{\vee}}
\newunicodechar{¬}{\ensuremath{\neg}}\newunicodechar{→}{\ensuremath{\to}}
\newunicodechar{↔}{\ensuremath{\leftrightarrow}}\newunicodechar{≤}{\ensuremath{\le}}
\newunicodechar{∈}{\ensuremath{\in}}\newunicodechar{⊢}{\ensuremath{\vdash}}
\newunicodechar{₂}{\ensuremath{_2}}\newunicodechar{₃}{\ensuremath{_3}}\newunicodechar{⁺}{\ensuremath{^+}}
% Lean listings are wider than the text font allows, so set them smaller.
\makeatletter\def\verbatim@font{\normalfont\ttfamily\footnotesize}\makeatother

% Shared by the PDF and the web build. Keep to simple \newcommand definitions.
% The short name of the main notion (alpha-plus-recursively Mahlo). Change it here and nowhere else.
% It is used both as an adjective ("the least \RM{} ordinal") and as a noun at the start of
% sentences, so it has to be an abbreviation or symbol, not a lower-case word.
\newcommand{\RM}{RM\textsuperscript{+}}
% Bold run-in head of a statement or front note: \head{Lemma 4.1}
\newcommand{\head}[1]{\textbf{#1.}}
% Named classes and predicates
\newcommand{\Mh}{\mathrm{Mh}}
\newcommand{\Sat}{\mathrm{Sat}}
\newcommand{\Mem}{\mathrm{Mem}}

\title{The least $\alpha^+$-recursively Mahlo ordinal is the least $\Pi_3$-reflecting ordinal}
% The date is stated here once; the PDF and the web page both take it from this line.
\newcommand{\paperversion}{2026-11-22}
\date{\paperversion}
\begin{document}
\maketitle

\begin{frontnote}\head{Abstract} Let $\Mh(\xi)$ be the $\xi$-th iterate of the recursively Mahlo operation along the
ordinals: $\Mh(\xi+1)$ is the class of ordinals that are $\Pi_2$-reflecting on $\Mh(\xi)$, with intersections at
limits. For an admissible ordinal $\alpha > \omega$ this iteration is defined at degree $\alpha$, but
$\alpha \notin \Mh(\alpha+1)$. We study higher degrees through
presentations: definable relations on $L_{\alpha}$, evaluated at smaller levels. Call $\alpha$
\emph{$\alpha^+$-recursively Mahlo} (\RM{}) if it lies in the Mahlo class of every node in the externally
well-founded part of each $\Sigma_1$ presentation at $\alpha$. We show that the least \RM{} ordinal is the least
$\Pi_3$-reflecting ordinal $K$, and that an admissible ordinal up to a certain bound is \RM{} exactly when it is
$\Pi_3$-reflecting or $\Pi_2$-reflecting on the $\Pi_3$-reflecting ordinals. Hence the first \RM{} ordinal that is not
$\Pi_3$-reflecting is the least ordinal $\gamma_1$ that is $\Pi_2$-reflecting on the $\Pi_3$-reflecting ordinals.
Restricting to $\Sigma_1$ linear orderings gives the same notion.
\end{frontnote}

\begin{frontnote}\head{Provenance} The proofs were found and written by an AI system. No human specialist has
reviewed them, and novelty is unconfirmed. Section 10 states what has been checked and how.
\end{frontnote}

\section{Introduction}

A cardinal $\kappa$ is greatly Mahlo if it is Mahlo of every degree below $\kappa^+$, the degrees at or above $\kappa$
being defined through canonical functions (Baumgartner, Taylor and Wagon). Below a weakly compact cardinal there are
stationarily many greatly Mahlo cardinals.

In the recursive setting, regular cardinals are replaced by admissible ordinals, stationary reflection by
$\Pi_2$-reflection on a class, and weak compactness by $\Pi_3$-reflection (Richter and Aczel). The ordinal analysis
of $\Pi_3$-reflection is due to Rathjen. Arai iterates the recursively Mahlo operation along a definable relation
that is evaluated in the current universe (Iterating, Definition 1.1), and shows that a $\Pi_3$-reflecting universe
lies in the class of every $n \in \omega$ up to which transfinite induction along a $\Sigma_1$ relation on $\omega$ holds
(Lemma 1.2; the same step is Lemma 3.33 of his Lectures, where Remark 3.34 notes that it is the only use of
$\Pi_3$-reflection). His main theorem is proof-theoretic: the theory of $\Pi_3$-reflecting universes is reducible to
the theory of limits of admissibles together with such iterations along one particular $\Sigma_1$ relation
(Theorem 1.3). Theorem A below is the same argument for $\Sigma_1$ presentations with parameters (Section 3.2), with
the well-founded part taken externally and with the strict reflection convention of Section 2. This note asks the
set-theoretic question that corresponds to Arai's reduction: which is the least ordinal in all these classes, and
which is the first one that is not $\Pi_3$-reflecting. We did not find Theorems B and C or Proposition D in the
sources we read; see Section 10.

Throughout, $\boldsymbol{A}$ is the class of $\Pi_3$-reflecting ordinals, $\boldsymbol{K}$ $= \min\, A$, and
$\boldsymbol{\gamma_1}$ is the least ordinal that is $\Pi_2$-reflecting on $A$. \RM{} is defined in 3.3.

\begin{statement}\head{Theorem A} Every $\Pi_3$-reflecting ordinal is \RM{}.
\end{statement}

\begin{statement}\head{Theorem B} No ordinal below $K$ is \RM{}. Hence the least \RM{} ordinal is $K$.
\end{statement}

\begin{statement}\head{Theorem C} An admissible ordinal $\gamma$ with $\omega < \gamma \le \gamma_1$ is \RM{} if and
only if $\gamma \in A$ or $\gamma$ is $\Pi_2$-reflecting on $A$. Hence the least \RM{} ordinal that is not
$\Pi_3$-reflecting is $\gamma_1$.
\end{statement}

\begin{statement}\head{Proposition D} For admissible $\alpha > \omega$, the definition of \RM{} is unchanged if only
$\Sigma_1$ presentations that define linear orderings of $L_{\alpha}$ are allowed.
\end{statement}

So \RM{} and $\Pi_3$-reflection have the same least element but are different classes.

\textbf{Overview of the proofs.} \emph{Why $\alpha^+$.} A node in the well-founded part of a $\Sigma_1$ relation on $L_{\alpha}$ has rank below
$\alpha^+$. The ranks that occur are exactly the order types of well-orderings that are $\Delta_1$-definable over
$L_{\alpha}$, and for every admissible $\alpha$ with $\omega < \alpha \le \gamma_1$ these fill all of
$\alpha^+$ (remark after Definition 3.3). \RM{} asks for membership in the Mahlo
class of every such node, as a greatly Mahlo cardinal $\kappa$ is Mahlo of every degree below $\kappa^+$. The class
of a node is not fixed by its rank: the relation is evaluated again at each smaller level (Definition 3.1, Examples
3.2).

\emph{The upper bound.} At an admissible level, membership in the class of a node is expressed by a $\Pi_3$ sentence
(Lemma 5.1). A $\Pi_3$-reflecting $\kappa$ therefore reflects ``$\varphi$, and the current level is in the class of
$c$'' for each predecessor $c$ of a node and each $\Pi_2$ sentence $\varphi$ true in $L_{\kappa}$, and induction on
the rank of the node gives Theorem A. This is Arai's argument.

\emph{The lower bound.} Fix a $\Pi_2$ sentence $\varphi_0(q_0)$. Section 6 defines a $\Sigma_1$ presentation $T$,
a tree whose nodes are finite increasing sequences of levels, each enumerated by $\omega$ in the next. Each step
handles one task: it supplies a witness for $\varphi_0(q_0)$, or it decides one $\Pi_3$ sentence
$\forall u\, \exists v\, \forall w\, \theta$ with constants already enumerated, either positively, by naming an
earlier limit level in which it holds, or negatively, by naming $u$ and promising $\forall v\, \exists w\, \neg\theta$;
later steps supply the witnesses $w$ that keep the promise. An infinite branch builds a $\Pi_3$-reflecting level
satisfying $\varphi_0(q_0)$ (Lemma 6.1). Below $K$, with $\varphi_0$ trivial, there is no such level, so the root is
in the well-founded part; ordinals that are not admissible are excluded by Example 3.2(1). To see that $\gamma$ is not in the class of the root, follow nodes whose promises are true
in $L_{\gamma}$. A promise is $\Pi_2$, so ``this node's promises are true'' is a $\Pi_2$ sentence that reflects down
together with the node, and the descent lemma (Lemma 4.1) turns this into non-membership in the class (Theorem M).
The only obstacle to extending such a node is a bad task: a $\Pi_3$ sentence true in $L_{\gamma}$ and in no earlier
limit level, which can be decided neither positively nor by a true promise. Such a sentence is false at every earlier limit level $L_b$, $b \ge \omega$,
containing its constant, so it cannot block a second time further down (Lemma 6.3). On a branch, the next record enumerates the
preceding selected level by $\omega$; Lemma 7.2 supplies these enumerations for every admissible $\gamma$ in the range
of Theorems B and C.

\emph{Theorem C.} The same tree is used with $\varphi_0(q_0)$ a $\Pi_2$ sentence that witnesses the failure of
$\Pi_2$-reflection on $A$ (Lemma 7.3). The converse direction applies Theorem A at the elements of $A$ to which
sentences are reflected (Lemma 7.4). Lemma 7.1 shows that $\gamma_1$ exists and Lemma 7.5 that it is not
$\Pi_3$-reflecting.

\emph{Proposition D.} A $\Sigma_1$ presentation is replaced by the Kleene--Brouwer ordering of its finite descending
sequences, recorded with edge witnesses; this does not need Sections 5--7.

\begin{center}\begin{tabular}{p{0.24\linewidth}p{0.70\linewidth}}
\toprule
Result & Uses \\
\midrule
Lemma 4.1 (descent) & Definition 3.1 only \\
Lemma 5.1 & Definition 3.1, (C1)--(C4), $\Sigma$-recursion and quantifier blocks (Section 2) \\
Theorem A & Lemma 5.1; Richter--Aczel, Theorem 1.8 \\
Lemma 6.1 & The definition of $T$ (Section 6.1), (C1), $\Delta_0$ absoluteness, quantifier blocks, dependent choice \\
Lemma 6.2 & The definitions of Sections 6.1 and 6.3, (C3) \\
Lemmas 6.3, 6.4 & The definition of a bad task (Section 6.3); for 6.4, the retention of constants (Quantifier blocks,
Section 2) \\
Theorem M & Lemmas 4.1 and 6.1--6.4, (C1), (C3), Definition 3.3 \\
Lemma 7.1 & (C1)--(C4); Richter--Aczel, Theorem 1.8; a countable elementary initial level of $L_{\omega_1}$ \\
Lemma 7.2 & (C1), (C2), (C6), $\Sigma$-recursion, the canonical well-ordering (Section 8), Lemma 7.1 \\
Theorem B & Theorems A and M, Lemmas 7.1 and 7.2, (C3), Example 3.2(1) \\
Lemma 7.3 & Theorem M, (C3) \\
Lemma 7.4 & Theorem A, (C5) \\
Lemma 7.5 & (C1)--(C4), Lemma 7.1, Step 4 of the proof of Lemma 5.1; Richter--Aczel, Theorem 1.8 \\
Theorem C & Theorem A, Example 3.2(1), Lemmas 7.1--7.5 \\
Corollary 7.6 & Lemmas 7.1 and 7.5, Theorem C \\
Proposition D & Definitions 3.1 and 3.3, (C1), (C2), $\Delta_0$ absoluteness, $\Sigma$-recursion, the order
certificates of Section 8, and the
well-foundedness of the Kleene--Brouwer ordering of a well-founded tree \\
\bottomrule
\end{tabular}\end{center}

\textbf{What to scrutinise.} \RM{} is a statement about \emph{presentations}: a definable relation is evaluated at
every smaller level, and its class depends on that evaluation and not only on its rank. \RM{} is therefore a
presentation-based iteration of recursive Mahloness. We do not prove that it corresponds to a recursive form of the
canonical-function or normal-ideal definition of greatly Mahlo cardinals, and the name ``recursively greatly Mahlo''
should be read with that reservation (Question 3 of Section 9). In particular Theorem B does not conflict with the
stationary-set theorem for cardinals quoted above: its proof enumerates each earlier level by $\omega$ inside the
ambient level, which is impossible at an inaccessible cardinal, and the cardinal notion evaluates degrees through
canonical functions modulo the nonstationary ideal, which is a different evaluation from Definition 3.1. The notion
is also sensitive to the complexity of the presentations (Section 3.4).

\section{Preliminaries}

\textbf{Background.} The following are standard and used without proof: $\Delta_0$ absoluteness; Kripke--Platek set
theory, of which every admissible level is a model, and its $\Sigma$-recursion theorem (Barwise, Theorem I.6.4); that the
map $\nu \mapsto L_{\nu}$ and the satisfaction relation of $L_{\nu}$ are $\Sigma_1$-definable, and so
$\Delta_1$-definable, over every admissible set (Devlin, Lemma I.9.10 and Section II.2); the condensation lemma (Devlin,
Theorem II.5.2); and Richter and Aczel's characterisation of the $\Pi_n$-reflecting ordinals, of which we use Theorem 1.8,
recalled below.

One point needs care. Formula complexity is counted below in the strict prenex sense, and no collection principle is
used to rearrange quantifiers, because the levels to which sentences are reflected need not be admissible. The usual
$\Sigma_1$ definitions of $L_{\nu}$ and of its satisfaction relation are justified by $\Sigma$-recursion, so at
admissible levels. We therefore
fix one explicit $\Sigma_1$ form of them: formulas asserting the existence of a \emph{certificate}, a set that records
the levels up to an ordinal with their satisfaction relations and is checked by a $\Delta_0$ formula. Certificates
are correct wherever they exist, and a fixed $\Pi_2$ sentence $G$ asserts that they exist for every ordinal; so they
can be used at every limit level satisfying $G$, not only at admissible ones. This is the standard construction of
$L$ inside Kripke--Platek set theory, set up so that the quantifier prefix of every formula used later can be read
off.

$L_{\alpha}$ is the $\alpha$-th level of the constructible hierarchy. Formulas are in the language of set theory with a
constant for each set; a sentence is evaluated in $L_{\beta}$ only if its constants are elements of $L_{\beta}$.

\textbf{Levels.} A \emph{level} is a set $L_{\nu}$; it is transitive, its ordinals are the ordinals below $\nu$, and
the earlier levels are those of smaller index. A set is \emph{countable in} $L_{\gamma}$ if it is empty or the image
of a surjection from $\omega$ that belongs to $L_{\gamma}$. A limit level, or an admissible level, is $L_{\nu}$ with $\nu$ a limit, or an admissible,
ordinal.

\textbf{Formula classes.} A formula is $\Delta_0$ if all its quantifiers are bounded. It is $\Pi_n$ if it is logically
equivalent to a formula $\forall x_1\, \exists x_2 \cdots \theta$ with $n$ alternating blocks of unbounded quantifiers
and $\theta \in \Delta_0$; $\Sigma_n$ is dual. This is the strict prenex convention of Richter and Aczel. No closure
under bounded quantification is assumed, and no collection principle is used to rearrange quantifiers.

\textbf{Reflection.} $\alpha$ is \emph{$\Pi_n$-reflecting} if every $\Pi_n$ sentence $\varphi(q)$ with
$q \in L_{\alpha}$ that is true in $L_{\alpha}$ is true in $L_{\beta}$ for some $\beta < \alpha$ with $q \in L_{\beta}$.
For a class $X$ of ordinals, $\mu$ is \emph{$\Pi_2$-reflecting on $X$} if every $\Pi_2$ sentence $\varphi(q)$ with
$q \in L_{\mu}$ that is true in $L_{\mu}$ is true in $L_{\beta}$ for some $\beta \in X \cap \mu$ with $q \in L_{\beta}$. A
limit ordinal $\alpha > \omega$ is $\Pi_2$-reflecting if and only if it is admissible (Richter--Aczel, Theorem 1.8).
Convention: the ordinals called $\Pi_n$-reflecting ($n \ge 2$) or $\Pi_2$-reflecting on a class are limit ordinals
above $\omega$. Tests include parameter-free sentences and sentences with any finite list of set parameters, each of
which must be an element of $L_{\beta}$ at the target $\beta$. The target need not be admissible.

\textbf{Absoluteness and recursion.} Two standard facts are used throughout (Barwise, Corollary I.8.6 and Theorem I.6.4). \emph{$\Delta_0$
absoluteness}: if $M \subseteq N$ are transitive sets, or $N$ is the universe, $\theta$ is $\Delta_0$ and
$x_1, \dots, x_n \in M$, then $M \models \theta(x_1, \dots, x_n)$ if and only if $N \models \theta(x_1, \dots, x_n)$.
A condition on given sets is called \emph{bounded} if it is expressed by a $\Delta_0$ formula with those sets as
constants; by absoluteness it does not matter in which transitive set containing them it is evaluated.
\emph{$\Sigma$-recursion}: an admissible level is a model of Kripke--Platek set theory, and there a
$\Sigma_1$-definable recursion step that gives a unique next set for each ordinal and each set-coded earlier history
defines a $\Sigma_1$-definable function by recursion on the ordinals, whose restriction to each ordinal of the level
is a set (the $\Sigma$-recursion theorem).

\textbf{Certificates.} Formulas without constants are coded by natural numbers, each code carrying a
declared finite context of free variables, and a sentence $\psi(q)$ is handled as the code of $\psi$ together with the
value $q$ of its one variable. The \emph{syntax data} are $\omega$ and the graphs of the primitive recursive functions
that decode codes, among them the marks for the bounded formulas and the prenex formulas $\forall x\, \exists y\, \theta$ and
$\forall u\, \exists v\, \forall w\, \theta$ with bounded $\theta$. Assignments are finite tuples, that is, functions
whose domain is a natural number. A \emph{certificate through $\xi$} is a
package $W$ of syntax data and a function on $\xi+1$ whose values are triples $W(\nu) = (M_\nu, F_\nu,
\mathrm{Tr}_\nu)$, the candidate level, assignment domain and truth table at $\nu$, subject to bounded conditions:
(V1) $F_\nu$ is a set of finite tuples from $M_\nu$ that contains the empty tuple and, with each of its members $s$ and
each $x \in M_\nu$, the extension of $s$ by $x$; (V2) $\mathrm{Tr}_\nu$ satisfies the Tarski clauses over $M_\nu$; (V3) $M_0 = \emptyset$, $M_\nu$ is the union of the earlier $M_\eta$ at limits, and
$M_{\nu+1}$ consists exactly of the subsets of $M_\nu$ defined by a formula code and an assignment from $F_\nu$
according to $\mathrm{Tr}_\nu$. Write $\mathrm{Cert}(W, \xi)$ for the conjunction of these conditions with those on
the syntax data; it is a $\Delta_0$ formula. The \emph{true} certificate through $\xi$ is the one in which, for every
$\nu \le \xi$, $M_\nu$ is $L_{\nu}$, $F_\nu$ is the set of all finite tuples from $L_{\nu}$ and $\mathrm{Tr}_\nu$ is
the satisfaction relation of $L_{\nu}$.

\begin{proof}[The definition in full]
Code the formulas without constants by natural numbers, in such a way that each
formula code carries a declared finite context of free variables, unused variables being allowed, and proper
subformulas have smaller codes than the formula. A sentence
$\psi(q)$ with a constant is handled as the code of $\psi$, with a one-variable context, together with the value $q$
of that variable. Fix primitive recursive functions that decide whether a number is a formula code and give its
constructor, the length of its context, its immediate subformulas, and the marks for the bounded formulas and for the
prenex formulas $\forall x\, \exists y\, \theta$ and $\forall u\, \exists v\, \forall w\, \theta$ with bounded
$\theta$. The \emph{syntax data} are $\omega$ together with the graphs of these functions and of the functions used in
their primitive recursive definitions. The conditions on syntax data are bounded and quantify only over the graphs
and over the natural numbers: that a set is $\omega$ says that it is a nonzero limit ordinal each of whose nonzero
members is a successor ordinal; and each graph is a set of tuples of the prescribed kind, is total and single-valued
on the prescribed finite power of $\omega$, and satisfies its defining equations at every argument in $\omega$. By
induction on the primitive recursive definitions and on their recursion arguments, only the true graphs satisfy these
conditions.

A \emph{certificate through $\xi$} is a package $W$ consisting of syntax data and a function with domain exactly
$\xi+1$, $\xi$ an ordinal, whose values
are triples $W(\nu) = (M_\nu, F_\nu, \mathrm{Tr}_\nu)$, which play the roles of the level, the assignment domain
and the truth table at $\nu$, such that the following bounded conditions hold. For finite tuples $s, t$ and a set
$x$, $\mathrm{Ext}(s, x, t)$ is the bounded condition that the domain of $t$ is that of $s$ plus one, $t$ agrees with
$s$ on the domain of $s$, and $t$ takes the value $x$ at the new place. All quantifiers in the conditions range over
$\xi+1$, $\omega$, the syntax data and the sets of the certificate.

\begin{itemize}
\item \textbf{(V1)} For each $\nu \le \xi$: every member of $F_\nu$ is a finite tuple from $M_\nu$, the empty tuple
  belongs to $F_\nu$, and
  $\forall s \in F_\nu\, \forall x \in M_\nu\, \exists t \in F_\nu\, \mathrm{Ext}(s, x, t)$.
\item \textbf{(V2)} For each $\nu \le \xi$: $\mathrm{Tr}_\nu$ is a set of pairs $(\varphi, s)$ of a formula code
  $\varphi \in \omega$ and a tuple $s \in F_\nu$ whose length is that of the context of $\varphi$. For every formula
  code $\varphi \in \omega$ and every $s \in F_\nu$ of that length, whether or not $(\varphi, s)$ belongs to
  $\mathrm{Tr}_\nu$, the Tarski biconditional for the constructor of $\varphi$ holds: for atomic formulas, negation
  and conjunction in the usual form, and for $\varphi = \exists x\, \varphi'$, with $x$ the last variable of the
  context of $\varphi'$, $(\varphi, s) \in \mathrm{Tr}_\nu$ if and only if there are $x \in M_\nu$ and $t \in F_\nu$
  with $\mathrm{Ext}(s, x, t)$ and $(\varphi', t) \in \mathrm{Tr}_\nu$; dually for $\forall$.
\item \textbf{(V3)} $M_0 = \emptyset$. For limit $\nu \le \xi$, every member of $M_\nu$ belongs to some $M_\eta$ with
  $\eta < \nu$, and every member of every such $M_\eta$ belongs to $M_\nu$. For $\nu < \xi$: for every formula code
  $\varphi \in \omega$ whose context has length $\ell + 1$ and every $s \in F_\nu$ of length $\ell$, some
  $z \in M_{\nu+1}$ consists exactly of the $x \in M_\nu$ such that $(\varphi, t) \in \mathrm{Tr}_\nu$ for the
  $t \in F_\nu$ with $\mathrm{Ext}(s, x, t)$; we say that $(\varphi, s)$ \emph{yields} $z$. Conversely, every member
  of $M_{\nu+1}$ is yielded by some such pair.
\end{itemize}
\end{proof}

\begin{itemize}
\item \textbf{(C1) Soundness.} If $\mathrm{Cert}(W, \xi)$ holds, then $W$ is the true certificate through $\xi$. For
  a formula code $\psi$ with one free variable and a statement such as ``$L_b \models \psi(q)$'', the existence of a
  certificate giving the positive verdict and the existence of one giving the negative verdict are the formulas
  \[ \exists W\, [\, \mathrm{Cert}(W, b) \wedge q \in M_b \wedge (\psi, \langle q \rangle) \in \mathrm{Tr}_b\, ] \]
  and the same formula with $\notin$ in place of the last $\in$. Both are $\Sigma_1$. Either verdict, when certified,
  is correct in any transitive ambient set, and in a transitive set that contains a certificate through $b$, with
  $q \in L_b$, the first formula is equivalent to the negation of the second. The truth of a $\Delta_0$ formula given
  by a code $\theta$ is treated in the same way, with a further quantifier $\exists \nu$ and the arguments required to
  lie in $M_\nu$.
  \begin{proof}[Self-contained proof of (C1)]
  The first statement is proved by induction on $\nu \le \xi$ (at $\nu = 0$ the only assignment is the empty one, which has the right length
  exactly for codes with empty context, and the definition $x = x$ without parameters yields $\emptyset$, so $M_1 = \{\emptyset\}$):
  $M_\nu = L_{\nu}$ by (V3) and the induction hypothesis for the earlier
  $M_\eta$ and $\mathrm{Tr}_\eta$; then $F_\nu$ is the set of all finite tuples from $L_{\nu}$ by (V1) and induction on
  length; then $\mathrm{Tr}_\nu$ is the satisfaction relation of $L_{\nu}$ by (V2) and induction on the construction of
  formulas. Since $\mathrm{Cert}$ is $\Delta_0$, a certificate that is an element of a transitive set and satisfies
  $\mathrm{Cert}$ there is the true one. So in any transitive ambient set the first formula implies
  $L_b \models \psi(q)$ and the second implies $L_b \nvDash \psi(q)$. For a $\Delta_0$ formula given by a code
  $\theta$ and arguments in $M_\nu$: by $\Delta_0$ absoluteness, $\theta$ is true of them if and only if it is true of
  them in $L_{\nu}$.
  \end{proof}
\item \textbf{(C2) Existence.} If $\alpha > \omega$ is admissible, every $\xi < \alpha$ has a certificate in
  $L_{\alpha}$. This is the standard construction of the levels and their satisfaction relations by
  $\Sigma$-recursion (Devlin, Section II.2).
  \begin{proof}[Self-contained proof of (C2)]
  Let $\alpha > \omega$ be admissible. Then $\omega \in L_{\alpha}$, and the graphs of the syntax data are obtained by
  $\Sigma$-recursion on $\omega$; they are elements of $L_{\alpha}$. $\Sigma$-recursion on $\nu$ then gives the
  successive triples $(M_\nu, F_\nu, \mathrm{Tr}_\nu)$, each step being a total $\Sigma_1$-definable operation in
  $L_{\alpha}$. At a limit stage, $M_\nu$ is the union of the earlier $M_\eta$. At a successor stage, for each formula
  code and each assignment in $F_\nu$ of compatible length, the subset of $M_\nu$ that the pair yields according to
  $\mathrm{Tr}_\nu$ is formed by $\Delta_0$ separation; these subsets are collected over the set of code--assignment
  pairs, which is an element of $L_{\alpha}$, and $M_{\nu+1}$ is the set of them, formed using collection and bounded
  separation. $F_\nu$ is the union over $n \in \omega$ of the sets of tuples of length $n$ from $M_\nu$, each obtained
  by recursion on $n$. $\mathrm{Tr}_\nu$ is built by recursion on the code, immediate subformulas having smaller codes:
  the entries for a code are separated, by its Tarski clause, from the set of its pairs with tuples in $F_\nu$, using
  the entries already obtained for smaller codes. The sequence through $\xi$, packaged with the syntax data, is the
  required certificate.
  \end{proof}
\item \textbf{(C3)} Let $\boldsymbol{G}$ be the sentence ``$\omega$ exists, every ordinal has a
  successor, and for every ordinal $\xi$ there is a certificate through $\xi$''. It is $\Pi_2$, and it holds in
  $L_{\alpha}$ for every admissible $\alpha > \omega$. If $L_{\beta} \models G$, then $\beta$ is a limit ordinal above
  $\omega$, and for every $\xi < \beta$ a certificate through $\xi$ is an element of $L_{\beta}$ and every such
  certificate is the true one. So a level satisfying $G$ is a limit level above $\omega$ in which
  truth in smaller levels and $\Delta_0$-truth can be evaluated: for $b < \beta$ and $q \in L_b$, exactly one of the
  two formulas of (C1) holds in $L_{\beta}$, and it gives the correct verdict; and every finite list of elements of
  $L_{\beta}$ lies in $L_{\nu}$ for some $\nu < \beta$. When a proof requires certificates to exist at a reflected
  target, the reflected sentence includes $G$, or the target is independently known to be admissible, or the
  sentence transferred to the target is an existential sentence that supplies the particular certificate (Lemmas 7.1
  and 7.2). So no bound on
  the height of certificates at non-admissible levels is needed.
  \begin{proof}[Self-contained proof of (C3)]
  The three conjuncts of $G$ have the forms
  $\exists x\, \theta_1$, $\forall \xi\, \exists y\, \theta_2$ and $\forall \xi\, \exists W\, \theta_3$ with bounded
  $\theta_1, \theta_2, \theta_3$ (for the third, ``$\xi$ is not an ordinal or $\mathrm{Cert}(W, \xi)$''), so $G$ is
  logically equivalent to a sentence $\forall \xi\, \forall \xi'\, \exists x\, \exists y\, \exists W\, \theta$ with
  bounded $\theta$: it is $\Pi_2$. (This prenex equivalence is the usual one for nonempty structures. The truth
  table of the empty structure $M_0$ is given separately by the clauses of (V2), with the empty assignment for
  sentences.) $G$ holds in $L_{\alpha}$ for every admissible $\alpha > \omega$, by (C2). If
  $L_{\beta} \models G$, then $\omega \in L_{\beta}$ and the ordinals of $L_{\beta}$ are closed under successor, so
  $\beta$ is a limit ordinal above $\omega$; and for every $\xi < \beta$ a certificate through $\xi$ is an element of
  $L_{\beta}$, and it is the true one by (C1).
  \end{proof}
\item \textbf{(C4)} There are fixed formulas $\Sat_2$ ($\Pi_2$) and $\Sat_3$ ($\Pi_3$) that define $\Pi_2$-truth and
  $\Pi_3$-truth of every level satisfying $G$. They range over prenex codes of the indicated form, invalid codes
  counting as false. For a bounded matrix and a finite tuple, an earlier level containing the tuple is chosen and its
  certified truth table is used, with the existential description of matrix truth in a final existential block and
  the universal description in a final universal block. Explicitly, $\Sat_2(\varphi, q)$ is
  \[ \forall x\, \exists y\, \exists \nu\, \exists W\, [\, \mathrm{Cert}(W, \nu) \wedge x, y, q \in M_\nu \wedge
     \mathrm{Val}_2\, ] \]
  and $\Sat_3(\psi, q)$ is
  \[ \forall u\, \exists v\, \forall w\, \forall \nu\, \forall W\, [\, \neg\mathrm{Cert}(W, \nu) \vee
     \neg(u, v, w, q \in M_\nu) \vee \mathrm{Val}_3\, ], \]
  where $\mathrm{Val}_2$ is the bounded condition ``by the syntax data of $W$, $\varphi$ is the code of a formula
  $\forall x\, \exists y\, \theta$ with $\theta$ bounded, and $(\theta, \langle q, x, y \rangle) \in \mathrm{Tr}_\nu$'',
  and $\mathrm{Val}_3$ is the same for $\psi$, $\forall u\, \exists v\, \forall w\, \theta$ and
  $\langle q, u, v, w \rangle$; each matrix is read in the context it inherits by (V2), the constant first and
  then the quantified variables in order. The
  brackets are bounded, so these are the strict prefixes, and no collection is used.
  \begin{proof}[Self-contained proof of (C4)]
  To see that they define truth, let $L_{\beta} \models G$. By (C3), any such arguments in $L_{\beta}$ lie in $L_{\nu}$
  for some $\nu < \beta$, a certificate through $\nu$ is an element of $L_{\beta}$, and every certificate in $L_{\beta}$
  is the true one; and by $\Delta_0$ absoluteness $\theta$ is true of them in $L_{\nu}$ if and only if it is true of
  them in $L_{\beta}$. So for a valid code, the bracket of $\Sat_2$ holds in $L_{\beta}$ for some $\nu, W$ exactly when
  $L_{\beta} \models \theta(x, y, q)$, and the bracket of $\Sat_3$ holds for all $\nu, W$ exactly when
  $L_{\beta} \models \theta(u, v, w, q)$. For an invalid code the first bracket fails for all $\nu, W$ and the second
  for some, so both formulas are false.
  \end{proof}
\item \textbf{(C5)} $\Sigma_1$ sentences persist upwards between levels, by $\Delta_0$ absoluteness.
\item \textbf{(C6) Condensation.} The transitive collapse of an elementary substructure of a limit level $L_{\eta}$ is a
  level (Devlin, Theorem II.5.2).
\end{itemize}

\textbf{Quantifier blocks.} In a limit level, a block of like quantifiers can be contracted to one quantifier over
finite tuples, using bounded quantifiers only: for a block of fixed length $\ell$, the universal quantifier over
tuples is guarded by the bounded predicate ``is a tuple of length $\ell$'', the existential quantifier over tuples
requires that predicate, and the entries are read by bounded quantifiers over the tuple. This is an equivalence in
every limit level, which contains the finite tuples of its elements ($L_{\omega}$ already does), but not in every
level: $L_1 = \{\emptyset\}$ contains no tuple of length $2$, so there $\exists x\, \exists y\, (x = y)$ is true and
its contraction is false. We write $\Pi_3$ sentences as
$\forall u\, \exists v\, \forall w\, \theta$, and we only ever use this contraction at limit levels. In the same way
a finite list of constants is replaced by one constant, the tuple of the list, which belongs to every limit level
containing its entries; this is why sentences are written $\varphi(q)$ with a single constant. A constant that a test
must keep at its target, whether or not it occurs in the sentence, is kept by conjoining its self-equality, as in
Lemmas 5.1 and 6.4.

\section{Presentations and the notion \RM{}}

\subsection{Iterating along the ordinals}

Let $\Mh(0)$ be the class of limit ordinals above $\omega$, $\Mh(\xi+1)$ the class of ordinals that are
$\Pi_2$-reflecting on $\Mh(\xi)$, and $\Mh(\lambda)$ the intersection of the earlier classes for limit $\lambda$. The
operation of $\Pi_2$-reflection on a class is monotone under inclusion, and its first iterate here is contained in
$\Mh(0)$; so by induction, taking intersections at limits, $\Mh(\zeta) \subseteq \Mh(\xi)$ for $\xi \le \zeta$. Then
$\Mh(1)$ is the class of admissible ordinals above $\omega$: by Richter--Aczel, Theorem 1.8, a limit
$\alpha > \omega$ is admissible if and only if it is $\Pi_2$-reflecting; reflection onto $\Mh(0)$ is in particular
$\Pi_2$-reflection; and conversely an admissible $\alpha > \omega$ reflects $\varphi(q) \wedge G$ to a level that is a
limit above $\omega$ by (C3). And $\Mh(2)$ is the class of recursively Mahlo ordinals. Here
$\Mh(\xi)$ depends only on the ordinal $\xi$, not on a relation that presents it.
This iteration is defined for all $\xi$, but no $\mu$ is in $\Mh(\mu+1)$: if $\mu$ were the least such ordinal,
it would be $\Pi_2$-reflecting on $\Mh(\mu)$, and reflecting a tautology would give $\beta < \mu$ in $\Mh(\mu) \subseteq \Mh(\beta+1)$. To express higher degrees
at $\mu$ one has to say how a degree is evaluated at smaller levels.

\subsection{Presentations and their classes}

A \emph{presentation} is a pair $P = (\rho, p)$ of a formula $\rho(x, y, z)$ and a set $p$. At a limit ordinal
$\mu > \omega$ with $p \in L_{\mu}$ it defines a relation on $L_{\mu}$: $c$ is a \emph{predecessor} of $a$ at $\mu$ if
$L_{\mu} \models \rho(c, a, p)$, and then $(c, a)$ is an \emph{edge} at $\mu$. Elements of $L_{\mu}$ are called
\emph{nodes}. $P$ is a $\Sigma_1$ presentation if
$\rho$ is $\Sigma_1$.

\begin{statement}\head{Definition 3.1} The \emph{class} $C(P, a)$ of a node $a$ is defined by recursion on $\mu$, for
all nodes simultaneously:

\begin{quote}
$\mu \in C(P, a)$ if and only if $\mu$ is a limit ordinal above $\omega$, $p$ and $a$ are elements of $L_{\mu}$, and
$\mu$ is $\Pi_2$-reflecting on $C(P, c)$ for every predecessor $c$ of $a$ at $\mu$.
\end{quote}
\end{statement}

The right-hand side refers only to $C(P, c) \cap \mu$. The relation is not required to be well-founded. A
\emph{descending sequence} from a node $a$ is a sequence $a = a_0, a_1, a_2, \dots$ in which each $a_{n+1}$ is a
predecessor of $a_n$. If there is no infinite descending sequence from $a$, the \emph{rank} of $a$ is the ordinal
defined by recursion as the supremum of the ordinals (rank of $c$)$\,+\,1$ over the predecessors $c$ of $a$. Both
notions refer to the relation at a given ordinal $\mu$.

\begin{statement}\head{Examples 3.2} (1) If $r$ is a well-founded relation that is a set and $\rho$ says
$(x, y) \in z$ with $p = r$, then $C(P, a)$ is $\Mh($rank of $a)$ restricted to the $\mu$ with $r, a \in L_{\mu}$. In
particular the relation $\{(0, 1)\}$ gives a node whose class is the class of admissible ordinals above $\omega$. (2)
Let $t = \{\{\emptyset\}\}$, which is not an ordinal, let $p = t$, and let $\rho(x, y, z)$ be the $\Delta_0$ formula
``$x$ is an ordinal, and either $y = z$ or $y$ is an ordinal with $x \in y$''. At every limit $\mu > \omega$, $t \in
L_{\mu}$ and the predecessors of $t$
are the ordinals below $\mu$, and each ordinal has the smaller ordinals as predecessors. As in (1), for an ordinal
$\xi$ the class $C(P, \xi)$ is $\Mh(\xi)$ restricted to the $\mu > \xi$, so $C(P, t)$ is the class of $\mu$ with
$\mu \in \Mh(\mu)$. The relation at $\mu$ is well-founded and $t$ has rank $\mu$ in it, so the rank of a node depends
on the ordinal at which the presentation is evaluated.
\end{statement}

In (1), by transitivity the predecessors of a node $a$ are the same at every $\mu$ with $r, a \in L_{\mu}$, and the
identity is proved by induction on the rank of $a$. A node without predecessors imposes only the domain conditions:
$\mu$ is a limit above $\omega$ with $r, a \in L_{\mu}$. For a predecessor $c$ of $a$, a test can be required to keep
$r$ and $c$ at its target by appending them as constants, so $\mu$ is $\Pi_2$-reflecting on $C(P, c)$ exactly when it
is $\Pi_2$-reflecting on $\Mh(\text{rank of } c)$, that is, when $\mu \in \Mh((\text{rank of } c) + 1)$. Over all
predecessors this gives the intersection of these classes, which is $\Mh(\text{rank of } a)$ because the iteration
is decreasing, at successor and at limit ranks alike. In (2) the same induction works with the predecessor ordinal
$\eta$ appended to each test, which forces the target above $\eta$; for $t$ the intersection over all $\eta < \mu$ is
the condition $\mu \in \Mh(\mu)$.

\subsection{The notion}

\begin{statement}\head{Definition 3.3} A limit ordinal $\alpha > \omega$ is \emph{$\alpha^+$-recursively Mahlo}
(\RM{}) if $\alpha \in C(P, a)$ for every $\Sigma_1$ presentation $P$ with $p \in L_{\alpha}$ and every node
$a \in L_{\alpha}$ in the well-founded part of the relation defined by $P$ at $\alpha$. Here the well-founded part of
a relation is the set of nodes below which it has no infinite descending sequence, taken in the universe and not
inside $L_{\alpha}$.
\end{statement}

By Example 3.2(1) every \RM{} ordinal is admissible. For admissible $\alpha$ let $\delta_{\alpha}$ be the supremum
of the order types of well-orderings with field contained in $L_{\alpha}$ that are $\Delta_1$-definable over
$L_{\alpha}$ with parameters. A node in the well-founded part at $\alpha$ has rank below $\delta_{\alpha}$: the ordering
$Q$ of Section 8, built from the node's histories, is such a well-ordering, and the rank of the node is at most the
order type of the histories that come before the empty history. Conversely, every $\xi < \delta_{\alpha}$ is the rank
of a node: a $\Delta_1$ well-ordering of type greater than $\xi$ is itself a $\Sigma_1$ presentation, and its element
at position $\xi$ has rank $\xi$. Always $\delta_{\alpha} \le \alpha^+$, the next admissible.

Let $\sigma$ be the least $\Sigma^1_1$-reflecting ordinal. By Gostanian's theorem (Gostanian, Section 2, application
(viii), and Corollary 3.2), for a countable admissible $\alpha$ such that every element of $L_{\alpha}$ is countable in $L_{\alpha}$, the
supremum of the order types of well-orderings of subsets of $\alpha$ that are $\Delta_1$-definable over $L_{\alpha}$
is $\alpha^+$ if $\alpha < \sigma$, and is smaller than $\alpha^+$ at $\alpha = \sigma$. That supremum is
$\delta_{\alpha}$: send each element of $L_{\alpha}$ to its position in the canonical well-ordering of $L$ (Section
8). Its predecessors there form a set in an earlier level, so $\Sigma$-recursion computes the position below $\alpha$,
and this map is a $\Delta_1$-definable injection of $L_{\alpha}$ into $\alpha$; it carries a $\Delta_1$ well-ordering
with field contained in $L_{\alpha}$ to one of the same type on a subset of $\alpha$. For admissible
$\omega < \alpha \le \gamma_1$, Lemma 7.1 gives external countability (its proof bounds $\gamma_1$ by a countable
ordinal), Lemma 7.2 gives local countability, and
Lemmas 7.1 and 7.5 give $\gamma_1 < \sigma$, since $\sigma$ is $\Pi_n$-reflecting for every $n$. Thus Gostanian's
theorem gives $\delta_{\alpha} = \alpha^+$ on this range. At other admissible $\alpha$ the definition measures degrees
below $\delta_{\alpha}$, which can be smaller than $\alpha^+$, as at $\sigma$. The rank of a node bounds what the node
presents; it does not determine its class.

\subsection{Persistence and the complexity of presentations}

By (C5), edges of a $\Sigma_1$ presentation at $\mu$ are edges at every $\gamma > \mu$. So a node in the well-founded
part at $\gamma$ is in the well-founded part at every smaller $\mu$ with the node in $L_{\mu}$. For a $\Pi_2$ presentation this
fails. Two statements about other complexity classes are not proved in this paper and are not used in it: that the
analogue of \RM{} for $\Pi_2$ presentations is equivalent to $\Pi_3$-reflection (a $\Pi_3$ sentence that fails to reflect can
be used to construct an obstruction of rank 2), and that Theorem A fails for $\Sigma_4$ presentations.

\section{A descent lemma}

\begin{statement}\head{Lemma 4.1} Let $P$ be a presentation and $Z$ a class of pairs $(\beta, c)$ with $\beta$ a limit
ordinal above $\omega$ and $c \in L_{\beta}$. Suppose that for every $(\beta, c) \in Z$ there are a predecessor $d$ of
$c$ at $\beta$ and a $\Pi_2$ sentence $\tau$, with constants in $L_{\beta}$ and true in $L_{\beta}$, such that
$(\beta', d) \in Z$ for every limit $\beta'$ with $\omega < \beta' < \beta$, $d$ and the constants of $\tau$ in
$L_{\beta'}$, and $L_{\beta'} \models \tau$. Then $\beta \notin C(P, c)$ for every $(\beta, c) \in Z$.
\end{statement}

\begin{proof}
Otherwise let $\beta$ be least such that $(\beta, c) \in Z$ and $\beta \in C(P, c)$ for some $c$, and take $d$ and
$\tau$. Since $d$ is a predecessor of $c$ at $\beta$, $\beta$ is $\Pi_2$-reflecting on $C(P, d)$. Applying this to
$\tau$ with $d$ adjoined as a constant gives $\beta' < \beta$ in $C(P, d)$ with $d$ and the constants of $\tau$ in $L_{\beta'}$
and $L_{\beta'} \models \tau$. Then $(\beta', d) \in Z$ and $\beta' \in C(P, d)$, contradicting the minimality of $\beta$.
\end{proof}

This lemma is formalised in Lean 4 with Mathlib (Appendix).

\section{Theorem A}

\begin{statement}\head{Lemma 5.1} For each $\Sigma_1$ formula $\rho$ there is a $\Pi_3$ formula $\Theta(p, a)$ such
that for every admissible $\alpha > \omega$ with $p, a \in L_{\alpha}$, $L_{\alpha} \models \Theta(p, a)$ if and only
if $\alpha \in C(P, a)$; and for every ordinal $\beta$ with $p, a \in L_{\beta}$, $L_{\beta} \models \Theta(p, a)$
implies $\beta \in C(P, a)$.
\end{statement}

\begin{proof}
The formula $\rho$ is fixed; the formulas $\Mem_\rho$ and $\Theta$ constructed below depend on it. Certificates and
their components $M_\nu$, $\mathrm{Tr}_\nu$ are as in Section 2.

\emph{Step 1: class certificates.} A \emph{class certificate through $\mu$} is a certificate through $\mu$ together
with sets $D(\nu) \subseteq M_\nu$, $\nu \le \mu$, such that for every $\nu \le \mu$ and every $x \in M_\nu$:
$x \in D(\nu)$ if and only if the clause of Definition 3.1 holds for $x$ at $\nu$ when membership of smaller ordinals
in classes is read from the $D(\eta)$, $\eta < \nu$. Written out, the clause is: $\nu$ is a limit ordinal above
$\omega$; $p, x \in M_\nu$; and for every $c \in M_\nu$ such that $\rho(c, x, p)$ is true according to
$\mathrm{Tr}_\nu$, every code $\varphi$ of a prenex $\Pi_2$ formula as in (C4) and every $q \in M_\nu$ such that
$\varphi(q)$ is true according to $\mathrm{Tr}_\nu$, there is $\eta < \nu$ with $q \in M_\eta$, $c \in D(\eta)$ and
$\varphi(q)$ true according to $\mathrm{Tr}_\eta$. Given the certificate, this condition is bounded, in both
directions: predecessors, codes of $\Pi_2$ sentences, their parameters and the candidate reflection ordinals all range
over sets of the certificate, and truth is read from its truth tables.

\emph{Step 2: soundness and existence.} Every class certificate has $D(\nu) = \{x \in L_{\nu} : \nu \in C(P, x)\}$ for
all $\nu \le \mu$. This is proved by induction on $\nu$. By (C1) the levels and truth tables of the certificate are the true
ones, so when the earlier $D(\eta)$ are as stated, the clause of Step 1 is the clause of Definition 3.1 at $\nu$,
except that it tests only prenex sentences with single quantifiers and one constant. That makes no difference,
because $\nu$ and all members of the classes are limit ordinals (Quantifier blocks, Section 2). In an admissible level
above $\omega$, class certificates through every $\mu$ exist, by $\Sigma$-recursion: $D(\nu)$ is obtained from
$W(\nu)$ and the earlier $D(\eta)$ by separating from $M_\nu$ with the bounded clause of Step 1.

\emph{Step 3: the formula.} Let $\Mem_\rho(\mu, p, c)$ be the $\Sigma_1$ formula ``there is a class certificate
through $\mu$ with $c \in D(\mu)$''. Let $\Theta(p, a)$ be the conjunction of $G$, $p = p \wedge a = a$, and
\[ \forall c\, \forall \varphi\, \forall q\, [\, \neg\rho(c, a, p) \vee \neg\Sat_2(\varphi, q) \vee \exists \mu\,
(\, p, c, q \in L_{\mu} \wedge \Mem_\rho(\mu, p, c) \wedge L_{\mu} \models \varphi(q)\, )\, ]. \]
In the last disjunct, $p, c, q \in L_{\mu}$ and $L_{\mu} \models \varphi(q)$ are read from the class certificate whose
existence $\Mem_\rho$ asserts, so this disjunct is $\Sigma_1$.

\emph{Step 4: $\Theta$ is $\Pi_3$.} Write $\rho$ as $\exists i\, \Phi_1$, $\Sat_2$ as
$\forall j\, \exists k\, \Phi_2$ as in (C4), with $k$ standing for a block, and the last disjunct as
$\exists h\, \Phi_3$, with $\Phi_1, \Phi_2, \Phi_3$ bounded. Then $\neg\rho$ is $\forall i\, \neg\Phi_1$ and
$\neg\Sat_2$ is $\exists j\, \forall k\, \neg\Phi_2$, so the displayed formula is logically equivalent to
$\forall c, \varphi, q, i\, \exists j, h\, \forall k\, [\neg \Phi_1 \vee \neg \Phi_2 \vee \Phi_3]$. Its conjunction
with the $\Pi_2$ sentence $G$ is again logically equivalent to a formula with this prefix, so $\Theta$ is $\Pi_3$ in
the strict sense.

\emph{Step 5: if $L_{\beta} \models \Theta(p, a)$, then $\beta \in C(P, a)$.} Let $L_{\beta} \models \Theta(p, a)$. Then
$L_{\beta} \models G$, so $\beta$ is a limit ordinal above $\omega$ by (C3). Let $c$ be a predecessor of $a$ at $\beta$
and $\varphi(q)$ a $\Pi_2$ sentence true in $L_{\beta}$, which we may take of the form $\forall x\, \exists y\, \theta$
with one constant. Then $L_{\beta} \models \Sat_2(\varphi, q)$ by (C4), so the last disjunct holds in $L_{\beta}$:
there are $\mu < \beta$ and a class certificate through $\mu$ in $L_{\beta}$ with $p, c, q \in L_{\mu}$,
$c \in D(\mu)$ and $L_{\mu} \models \varphi(q)$. By Step 2 this successful class certificate gives actual membership,
$\mu \in C(P, c)$. So $\beta$ is $\Pi_2$-reflecting on $C(P, c)$ for every predecessor $c$ of $a$, that is,
$\beta \in C(P, a)$.

\emph{Step 6: the converse at admissible levels.} Let $\alpha > \omega$ be admissible with $\alpha \in C(P, a)$. Then
$L_{\alpha} \models G$ by (C3). Let $c, \varphi, q \in L_{\alpha}$ with $L_{\alpha} \models \rho(c, a, p)$ and
$L_{\alpha} \models \Sat_2(\varphi, q)$. Then $c$ is a predecessor of $a$ at $\alpha$ and $\varphi(q)$ is a $\Pi_2$
sentence true in $L_{\alpha}$, so there is $\mu < \alpha$ in $C(P, c)$ with $q \in L_{\mu}$ and
$L_{\mu} \models \varphi(q)$; also $p, c \in L_{\mu}$ by Definition 3.1. By Step 2 a class certificate through $\mu$ is
an element of $L_{\alpha}$ and has $c \in D(\mu)$. So the last disjunct holds in $L_{\alpha}$.
\end{proof}

\begin{statement}\head{Theorem A} Let $\kappa$ be $\Pi_3$-reflecting, $P$ a $\Sigma_1$ presentation with
$p \in L_{\kappa}$, and $a$ a node in the well-founded part of the relation defined by $P$ at $\kappa$. Then
$\kappa \in C(P, a)$. Hence $\kappa$ is \RM{}.
\end{statement}

\begin{proof}
$\kappa$ is admissible and above $\omega$. Argue by induction on the rank of $a$. Let $c$ be a predecessor of $a$ at
$\kappa$. By the induction hypothesis $\kappa \in C(P, c)$, so $L_{\kappa} \models \Theta(p, c)$. Let $\varphi(q)$ be
a $\Pi_2$ sentence true in $L_{\kappa}$. Then $\varphi(q) \wedge \Theta(p, c)$ is a $\Pi_3$ sentence true in
$L_{\kappa}$, so it holds in $L_{\beta}$ for some $\beta < \kappa$ with $p, c, q \in L_{\beta}$. By Lemma 5.1,
$\beta \in C(P, c)$. So $\kappa$ is $\Pi_2$-reflecting on $C(P, c)$ for every predecessor $c$ of $a$, that is,
$\kappa \in C(P, a)$.
\end{proof}

The induction is carried out in the universe, along the well-founded part. This is the argument of Arai for
iterations of the recursively Mahlo operation along relations interpreted in the current level.

\section{A tree that searches for a $\Pi_3$-reflecting level}

Fix a $\Pi_2$ sentence $\varphi_0(q_0) = \forall x\, \exists y\, \theta_0(x, y, q_0)$ with $\theta_0 \in \Delta_0$. We
define a $\Sigma_1$ presentation $T$ whose infinite branches produce a level that is $\Pi_3$-reflecting and satisfies
$\varphi_0(q_0)$.

\subsection{Definition of $T$}

A \emph{node of depth $k$} is a finite sequence of \emph{records} numbered $0, \dots, k$, where:

\begin{itemize}
\item record $j$ contains a limit ordinal $a_j > \omega$, its \emph{stage}, with $q_0 \in L_{a_0}$ and
  $a_0 < a_1 < \cdots < a_k$;
\item for $j \ge 1$, record $j$ also contains a surjection $f_{j-1}$ from $\omega$ onto $L_{a_{j-1}}$ that is an element
  of $L_{a_j}$, and an element $c_j$ of $L_{a_j}$, the \emph{answer} to task $j-1$.
\end{itemize}

The \emph{last stage} of the node is $a_k$. A \emph{child} of a node is an
extension of it by one record, the new record; for a node of depth $k$ this is record $k+1$, with stage
$a_{k+1}$.

Fix a primitive recursive enumeration of \emph{tasks}, with task 0 empty, in which every task below occurs and task
$k$ refers only to surjections $f_i$ with $i < k$ and answers $c_j$ with $1 \le j \le k$. The next task of a node of
depth $k$ is task $k$. So task $k$ mentions only surjections and answers already present in a node of depth $k$,
and it is answered by $c_{k+1}$, the answer in the new record. An answer is read according to the task it
answers. There are three kinds of nonempty task: $\mathrm{Wit}(i, n)$ asks for a witness for $\varphi_0$,
$\mathrm{Refl}(i, n)$ asks to decide one instance of $\Pi_3$-reflection, and $\mathrm{Keep}(j, i, n)$ asks to keep, in
one instance, the promise made by a negative answer.

\begin{itemize}
\item \textbf{$\mathrm{Wit}(i, n)$.} Let $x = f_i(n)$. The answer is some $y$ with $\theta_0(x, y, q_0)$.
\item \textbf{$\mathrm{Refl}(i, n)$.} If $f_i(n)$ is a pair $(\psi, q)$ with $\psi$ a $\Pi_3$ formula
  $\forall u\, \exists v\, \forall w\, \theta$, the answer is either
  \begin{itemize}
  \item $(+, b)$, where $b$ is a \textbf{limit} ordinal with $\omega \le b$ below the stage of the new record,
    $q \in L_b$, and $L_b \models \psi(q)$; or
  \item $(-, u)$, where $u$ is any set. This is a promise that $\forall v\, \exists w\, \neg\theta(u, v, w, q)$.
    Nothing is verified when it is made.
  \end{itemize}
  We then call the task a \emph{Refl-task for $(\psi, q)$}.
\item \textbf{$\mathrm{Keep}(j, i, n)$.} If $c_j$ is an answer $(-, u)$ to a Refl-task for $(\psi, q)$ and
  $v = f_i(n)$, the answer is some $w$ with $\neg\theta(u, v, w, q)$.
\item Every other task is empty and is answered by $\emptyset$.
\end{itemize}

\textbf{Example.} Suppose, for illustration only, that the enumeration begins with task 1 $= \mathrm{Wit}(0, 0)$,
task 2 $= \mathrm{Refl}(0, n)$ and task 3 $= \mathrm{Keep}(3, 1, m)$ for some numbers $n, m$. A node of depth 3
consists of four records.
Record 0 is a stage $a_0$ with $q_0 \in L_{a_0}$. Record 1 is $(a_1, f_0, c_1)$ with $f_0$ a surjection from $\omega$
onto $L_{a_0}$ and $c_1 = \emptyset$, the answer to the empty task 0. Record 2 is $(a_2, f_1, c_2)$, where $c_2$
answers task 1: it is some $y \in L_{a_2}$ with $\theta_0(f_0(0), y, q_0)$. Record 3 is $(a_3, f_2, c_3)$, where $c_3$
answers task 2. If $f_0(n)$ is a pair $(\psi, q)$ of the required form, then $c_3$ is either $(+, b)$ for a limit
ordinal $b$ with $\omega \le b < a_3$, $q \in L_b$ and $L_b \models \psi(q)$, or $(-, u)$ for some $u \in L_{a_3}$;
otherwise task 2 is empty and $c_3 = \emptyset$. The next task of this node is task 3, which refers to the answer
$c_3$ and the surjection $f_1$. If $c_3 = (-, u)$, a child adds record 4, which is $(a_4, f_3, c_4)$ with $f_3$ a surjection from $\omega$ onto
$L_{a_3}$ and $c_4 = w$ an answer with $\neg\theta(u, f_1(m), w, q)$. This is one instance of the promise made by $c_3$; the other instances are the tasks
$\mathrm{Keep}(3, i, n')$, which occur later in the enumeration. If $c_3$ is not of the form $(-, u)$, task 3 is empty
and $c_4 = \emptyset$.

The predecessors of a node are its children, and a \emph{root} $r$ has the nodes of depth 0 as predecessors. Take
$r = \emptyset$, the empty sequence: every node has positive length as a sequence of records, and $r$ belongs to
every level considered. ``$s$
is a node'' and ``$d$ is a child of $s$'' are expressed by $\Sigma_1$ formulas with the constant $q_0$: they assert
the existence of a certificate through the last stage, of $s$ or of $d$ respectively, according to which the bounded
conditions above hold. Given the certificate the conditions are indeed bounded: the levels $L_{a_j}$ and $L_b$ are
given by the certificate, the instances of $\theta_0$ and $\theta$ are read from the truth table at the last stage $a_k$, whose level $L_{a_k}$
contains their arguments, $L_b \models \psi(q)$ is read from the truth table at $b$, and the enumeration of tasks, being
primitive recursive, is definable over $L_{\omega}$ and is read from the truth table at $\omega$. By (C1) any such
certificate gives a true edge, so $T$ is a $\Sigma_1$ presentation.

\subsection{Branches}

\begin{statement}\head{Lemma 6.1} If the root is not in the well-founded part of the relation defined by $T$ at $\mu$,
there is a limit ordinal $\lambda \le \mu$ with $q_0 \in L_{\lambda}$ such that $\lambda$ is $\Pi_3$-reflecting and
$L_{\lambda} \models \varphi_0(q_0)$.
\end{statement}

\begin{proof}
Take a sequence of nodes, each a child of the previous one, with stages $a_0 < a_1 < \cdots$ in $L_{\mu}$, and let
$\lambda$ be their supremum. Then $L_{\lambda}$ is the union of the ranges of the $f_j$. Every $x \in L_{\lambda}$ is
$f_i(n)$ for some $i, n$, and task $\mathrm{Wit}(i, n)$ was answered by a $y$ with $\theta_0(x, y, q_0)$; hence
$L_{\lambda} \models \varphi_0(q_0)$. Let $\psi(q) = \forall u\, \exists v\, \forall w\, \theta$ be a $\Pi_3$ sentence
with $q \in L_{\lambda}$. The pair $(\psi, q)$ is $f_i(n)$ for some $i, n$, and task $\mathrm{Refl}(i, n)$ was
answered. If the
answer was $(+, b)$, then $b$ is a limit below $\lambda$ with $q \in L_b$ and $L_b \models \psi(q)$. If it was
$(-, u)$, then every $v \in L_{\lambda}$ is $f_{i'}(n')$ for some $i', n'$, and the corresponding Keep-task was answered by
a $w$ with $\neg\theta(u, v, w, q)$; so $L_{\lambda} \models \forall v\, \exists w\, \neg\theta(u, v, w, q)$ and
$\psi(q)$ is false in $L_{\lambda}$. Hence every sentence of the form $\forall u\, \exists v\, \forall w\, \theta$ true
in $L_{\lambda}$ is true in an earlier limit level containing its constant. A $\Pi_3$ sentence with quantifier blocks is
equivalent to one of that form in every limit level, so $\lambda$ is $\Pi_3$-reflecting.
\end{proof}

\subsection{Correct nodes}

Let $\beta$ be a limit ordinal above $\omega$ and $s$ a node that is an element of $L_{\beta}$.

\begin{itemize}
\item $s$ is \emph{correct at $\beta$} if every promise in $s$ is true in $L_{\beta}$: for every answer $(-, u)$ in
  $s$ to a Refl-task for $(\psi, q)$, where $\theta$ is the matrix of $\psi$,
  $L_{\beta} \models \forall v\, \exists w\, \neg\theta(u, v, w, q)$. (Answers to other tasks are not promises, even
  if they happen to have the form $(-, u)$.) Let $\mathrm{corr}(s)$ be the conjunction of these
  finitely many $\Pi_2$ sentences, each with its own formula $\theta$ and constants $u, q$.
\item A Refl-task for $(\psi, q)$ is \emph{bad at $\beta$} if $L_{\beta} \models \psi(q)$ and $L_b \nvDash \psi(q)$ for
  every limit $b$ with $\omega \le b < \beta$ and $q \in L_b$. Tasks of the other kinds are never bad.
\item $\boldsymbol{\chi}$ is the $\Pi_2$ sentence $\varphi_0(q_0) \wedge G \wedge$ ``every ordinal is smaller than some
  limit ordinal'' $\wedge$ ``every $L_a$ is countable'', that is, for every ordinal $a$, either $L_a$ is empty or
  there is a surjection from $\omega$ onto $L_a$, the level $L_a$ being supplied by a certificate.
\end{itemize}

If $L_{\beta} \models \chi$, then $\beta$ is a limit of limit ordinals above $\omega$. Every admissible level above
$\omega$ satisfies $G$ and the clause ``every ordinal is smaller than some limit ordinal''.

\begin{statement}\head{Lemma 6.2} Suppose $L_{\beta} \models \chi$, the node $s \in L_{\beta}$ is correct at $\beta$,
and the next task of $s$ is not a Refl-task that is bad at $\beta$. Then $s$ has a child in $L_{\beta}$ that is correct
at $\beta$ and is a predecessor of $s$ at $\beta$.
\end{statement}

\begin{proof}
Choose the answer to the next task. For a Wit-task, $y$ exists because $\varphi_0(q_0)$ holds in $L_{\beta}$. For a
Keep-task, $w$ exists because the promise $(-, u)$ is true in $L_{\beta}$. For a Refl-task for $(\psi, q)$ that is not
bad at $\beta$: if $\psi(q)$ is false in $L_{\beta}$, some $u \in L_{\beta}$ has
$L_{\beta} \models \forall v\, \exists w\, \neg\theta(u, v, w, q)$, and the answer $(-, u)$ is correct at $\beta$; if
$\psi(q)$ is true in $L_{\beta}$, there is a limit $b < \beta$ with $q \in L_b$ and $L_b \models \psi(q)$, and
$(+, b)$ is allowed. By the last clause of $\chi$ there is a surjection $f \in L_{\beta}$ from $\omega$ onto $L_{a_k}$, where $a_k$ is the last stage of $s$.
The answer, $f$, and $b$ if present are finitely many elements of $L_{\beta}$, and $\beta$ is a limit of limits, so
there is a limit ordinal $a < \beta$ above the last stage with all of them in $L_a$. The extension of $s$ by the
record $(a, f,$ answer$)$ is an element of $L_{\beta}$, and by $G$ there is a certificate through $a$ in $L_{\beta}$,
which makes the $\Sigma_1$ edge formula true in $L_{\beta}$.
\end{proof}

\begin{statement}\head{Lemma 6.3} If a Refl-task for $(\psi, q)$ is bad at $\beta$, it is not bad at any limit
$\beta'$ with $\omega < \beta' < \beta$ and $q \in L_{\beta'}$.
\end{statement}

\begin{proof}
By the definition of bad at $\beta$, $\psi(q)$ is false in $L_{\beta'}$.
\end{proof}

Badness, and so Lemma 6.3, concerns earlier limit levels only; a sentence whose Refl-task is bad at $\beta$ can be
true at a smaller successor level. For example, take $\varphi_0$ to be $\forall x\, \exists y\, (y = x)$ with $q_0 =
\emptyset$, write $\chi$ with its blocks contracted as $\forall x\, \exists y\, \theta_\chi(x, y, q_0)$, and let
$\psi(q_0)$ be $\forall u\, \exists v\, \forall w\, [\, \theta_\chi(u, v, q_0) \vee (v$ is an ordinal $\wedge\, (w$ is
an ordinal $\to w \in v \vee w = v))\, ]$, which allows the witness $v$ alternatively to be a largest ordinal. This is
strict $\Pi_3$. At infinite limit levels, which have no largest ordinal, $\psi(q_0)$ is equivalent to $\chi$; and
$\psi(q_0)$ holds in $L_{\omega+1}$, whose largest ordinal is $\omega$. So the Refl-task for $(\psi, q_0)$ is bad at
the least limit at which $\chi$ holds, although $\psi(q_0)$ is true in $L_{\omega+1}$.

\begin{statement}\head{Lemma 6.4} If $L_{\beta} \models \chi$ and a task is not bad at $\beta$, there is a $\Pi_2$
sentence $\varsigma$ with constants in $L_{\beta}$, true in $L_{\beta}$, such that the task is not bad at any limit
$\beta'$ with $\omega < \beta' < \beta$, the constants of $\varsigma$ in $L_{\beta'}$, and $L_{\beta'} \models \varsigma$.
\end{statement}

\begin{proof}
For tasks that are not Refl-tasks take $\varsigma$ trivial. For a Refl-task for $(\psi, q)$ with $\psi(q)$ false in
$L_{\beta}$, take $u$ with $L_{\beta} \models \forall v\, \exists w\, \neg\theta(u, v, w, q)$ and let $\varsigma$ be
$(u = u) \wedge (q = q) \wedge \forall v\, \exists w\, \neg\theta(u, v, w, q)$; then $u, q \in L_{\beta'}$ and
$\psi(q)$ is false in $L_{\beta'}$. For a Refl-task with $\psi(q)$ true in $L_{\beta}$, take a limit
$b < \beta$ with $q \in L_b$ and $L_b \models \psi(q)$, and let $\varsigma$ be $b = b$; then $b < \beta'$, so
$q \in L_b \subseteq L_{\beta'}$, and the task is not bad at $\beta'$.
\end{proof}

\subsection{The theorem}

Theorem M contains two statements, proved by two different arguments: that $\beta$ is not in the class of the root
(by descent, Lemma 4.1) and that the root is in the well-founded part at $\beta$ (no infinite branch, Lemma 6.1). For
the first, Lemma 4.1 is applied to the family $Z$ of the following pairs, where in both cases $q_0 \in L_{\beta}$ and
$L_{\beta} \models \chi$:

\begin{itemize}
\item the pairs $(\beta, r)$ with the root $r$;
\item the pairs $(\beta, s)$, where $s \in L_{\beta}$ is a node that is correct at $\beta$ and whose next
  task is not bad at $\beta$.
\end{itemize}

For each pair $(\beta, s)$ in $Z$ the proof gives a predecessor $d$ of $s$ at $\beta$, which is a node correct at
$\beta$, and a $\Pi_2$ sentence $\tau$ true in $L_{\beta}$, and shows that $(\beta', d)$ is in $Z$ for
every limit $\beta'$ with $\omega < \beta' < \beta$, $d$ and the constants of $\tau$ in $L_{\beta'}$, and
$L_{\beta'} \models \tau$. The pair
$(\beta, d)$ itself need not be in $Z$, because the next task of $d$ may be bad at $\beta$; Lemma 4.1 does not require
it. The only obstacle to a correct node having a correct child is a $\Pi_3$ sentence that is true in $L_{\beta}$ and in no
earlier limit level, and by Lemma 6.3 such a sentence cannot be an obstacle again further down.

\begin{statement}\head{Theorem M} Let $\beta$ be an ordinal with $q_0 \in L_{\beta}$ and $L_{\beta} \models \chi$. Then
$\beta \notin C(T, r)$. If moreover there is no $\Pi_3$-reflecting $\lambda \le \beta$ with $q_0 \in L_{\lambda}$ and
$L_{\lambda} \models \varphi_0(q_0)$, then the root is in the well-founded part of the relation defined by $T$ at
$\beta$, and so $\beta$ is not \RM{}.
\end{statement}

\begin{proof}
Let $Z$ be the family of pairs defined before the theorem. Every $\beta$ occurring in it is a
limit of limit ordinals above $\omega$, because $L_{\beta} \models \chi$. We verify the hypothesis of Lemma 4.1 for a
pair $(\beta, s) \in Z$, where $s$ is the root or a node.

\emph{The predecessor.} If $s$ is a node, let $d$ be a child of $s$ correct at $\beta$ (Lemma 6.2). If $s = r$, let
$d$ be a node of depth 0 in $L_{\beta}$; one exists because $\beta$ is a limit of limits, and it has no answers,
so it is correct at $\beta$. In both cases $d$ is a predecessor of $s$ at $\beta$: for a child this is part of Lemma
6.2, and for a node of depth 0 the certificate required by the edge formula is supplied by $G$.

\emph{The sentence.} Let $\tau$ be $\chi \wedge \mathrm{corr}(d) \wedge \varsigma$, with the constant $d$ adjoined,
where $\varsigma$ is given by Lemma 6.4 for the next task of $d$ if that task is not bad at $\beta$ and is trivial
otherwise. It is $\Pi_2$ and true in $L_{\beta}$. Note that $\tau$ never asserts that a task which is bad at $\beta$ is
not bad.

\emph{Passing down.} Let $\beta'$ be a limit with $\omega < \beta' < \beta$, $d$ and the constants of $\tau$ in $L_{\beta'}$, and
$L_{\beta'} \models \tau$. Then $L_{\beta'} \models \chi$; $d$ is a node there (its last stage is below $\beta'$ because
$d \in L_{\beta'}$; by $G$ a certificate through that stage belongs to $L_{\beta'}$; its tables are the true ones, so the
bounded checks for the record $d$ succeed); and $d$ is correct at $\beta'$, because $\mathrm{corr}(d)$ holds in
$L_{\beta'}$. It remains to see that the next task of $d$ is not bad at $\beta'$. If it was not bad at $\beta$, this
is Lemma 6.4, since the constants of $\varsigma$ are in $L_{\beta'}$ and $L_{\beta'} \models \varsigma$. If it was bad at
$\beta$, it is a Refl-task for a pair $(\psi, q) = f_i(n)$, where $f_i$ is one of the surjections in $d$. Since
$d \in L_{\beta'}$ and $L_{\beta'}$ is transitive, $f_i$, the pair $(\psi, q)$ and $q$ are elements of $L_{\beta'}$, so
Lemma 6.3 applies. So $(\beta', d) \in Z$.

Lemma 4.1 gives $\beta \notin C(T, s)$ for all $(\beta, s) \in Z$, in particular $\beta \notin C(T, r)$. For the
second statement, if the root were not in the well-founded part at $\beta$, Lemma 6.1 would give a $\lambda$ excluded
by the hypothesis; so the root is in the well-founded part at $\beta$ and $\beta \notin C(T, r)$, and $\beta$ is not
\RM{} by Definition 3.3.
\end{proof}

\section{Theorems B and C}

\begin{statement}\head{Lemma 7.1} Every $\Pi_4$-reflecting ordinal is $\Pi_2$-reflecting on $A$ and belongs to $A$.
Such ordinals exist; in particular $\gamma_1$ exists.
\end{statement}

\begin{proof}
Let $R_3$ be the conjunction of $G$ with $\forall\psi\, \forall q\, [\, \neg\Sat_3(\psi, q) \vee \exists b\, (\, b$ is
a limit ordinal above $\omega \wedge q \in L_b \wedge L_b \models \psi(q)\, )\, ]$, the last clause expressed through
a certificate and $\psi$ ranging over sentences $\forall u\, \exists v\, \forall w\, \theta$. It is $\Pi_4$. A level
satisfying it reflects every such sentence into a limit level, where block contraction is valid, so it is
$\Pi_3$-reflecting, by (C1) and (C4). A $\Pi_4$-reflecting $\delta$ is $\Pi_3$-reflecting and admissible above
$\omega$; reflecting $\psi(q) \wedge G$ for each true $\psi(q)$ gives limit targets above $\omega$, so
$L_{\delta} \models R_3$. Given a $\Pi_2$ sentence $\varphi(q)$ true in $L_{\delta}$, reflect $\varphi(q) \wedge R_3$.
For existence, take a countable $\alpha$ with $L_{\alpha}$ a proper elementary substructure of $L_{\omega_1}$. Such
$\alpha$ exist by the standard L\"owenheim--Skolem argument (Devlin, Sections II.5--6): the
chain $(L_{\xi})_{\xi < \omega_1}$ has countable stages, and closing under the countably many Skolem functions for
$L_{\omega_1}$ gives such an $\alpha < \omega_1$. For each sentence $\varphi(q)$ true in $L_{\alpha}$, $L_{\omega_1}$
satisfies the existential sentence ``there are a limit ordinal $b > \omega$ and a certificate through $b$ with
$q \in M_b$ and $\varphi(q)$ true according to its truth table at $b$'': $b = \alpha$ and its certificate from (C2)
witness it ($L_{\omega_1}$ is admissible, $\omega_1$ being a regular uncountable cardinal). By elementarity the
sentence holds in $L_{\alpha}$, and by (C1) the $b < \alpha$ it supplies is a limit above $\omega$ with
$q \in L_b$ and $L_b \models \varphi(q)$.
\end{proof}

The next lemma is needed because Lemma 6.2 extends a node with last stage $a_k$ by a surjection from $\omega$ onto
$L_{a_k}$, and that surjection comes from the last clause of $\chi$. It shows that this clause fails in an admissible
$L_{\gamma}$ only if some $\xi_0 < \gamma$ is $\Pi_n$-reflecting for every $n$, so it holds in the range of Theorems B
and C. The proof shows by
condensation that the first level which $L_{\gamma}$ cannot enumerate is a $\Sigma_1$-elementary substructure of
$L_{\gamma}$. This is the analogue inside $L_{\gamma}$ of the standard fact that $L_{\omega_1}$ is a
$\Sigma_1$-elementary substructure of $L$, proved by the same L\"owenheim--Skolem and condensation argument (Devlin,
Lemmas II.5.10 and II.5.11).

\begin{statement}\head{Lemma 7.2} Let $\gamma > \omega$ be admissible and suppose some $L_a$ with $a < \gamma$ is not countable in $L_{\gamma}$.
Then some $\xi_0 < \gamma$ is $\Pi_n$-reflecting for every $n$. Hence, for every admissible $\gamma$ with $\omega < \gamma < K$ and for every admissible $\gamma$ with
$\omega < \gamma \le \gamma_1$, every $L_a$ with $a < \gamma$ is countable in $L_{\gamma}$.
\end{statement}

\begin{proof}
In $L_{\gamma}$ every finite level and $L_{\omega}$ are countable. Let $\xi_0$ be the least $a$ such that $L_a$ is
nonempty and $L_{\gamma}$ contains no surjection from $\omega$ onto $L_a$. If a nonempty $L_a$ is countable in
$L_{\gamma}$, then its enumeration, formula codes and finite parameter tuples enumerate $L_{a+1}$ in $L_{\gamma}$, using
satisfaction for the set $L_a$. So $\xi_0$ is a limit above $\omega$, and $L_{\bar{\xi}}$ is countable in $L_{\gamma}$
for every $\bar{\xi} < \xi_0$.

\emph{$L_{\xi_0}$ is a $\Sigma_1$-elementary substructure of $L_{\gamma}$.} A countable elementary substructure of a level containing the
parameter and a witness, formed inside $L_{\gamma}$, collapses by condensation to a level that is countable in
$L_{\gamma}$, hence to a level below $\xi_0$, which contains a witness. This is the standard argument referred to
before the lemma; it is written out in the self-contained proof below.

\emph{Reflection.} Now let $\varphi(q)$ be true in $L_{\xi_0}$. Then $L_{\gamma}$ satisfies the $\Sigma_1$ sentence
``some limit level of height above $\omega$ containing $q$ satisfies $\varphi(q)$'', witnessed by $\xi_0$ through a
certificate (the condition on the height is bounded in the certificate). So $L_{\xi_0}$ satisfies it, giving such a
level below $\xi_0$.
Thus $\xi_0$ is $\Pi_n$-reflecting for every $n$. In particular $K \le \xi_0 < \gamma$, and by Lemma 7.1 $\xi_0$ is
$\Pi_2$-reflecting on $A$, so $\gamma_1 \le \xi_0 < \gamma$.
\end{proof}

\begin{proof}[Self-contained proof of the elementary substructure step]
Let $x \in L_{\bar{\xi}}$ with $\omega \le \bar{\xi} < \xi_0$, and suppose $L_{\gamma} \models \exists z\, \theta(z, x)$
with $\theta \in \Delta_0$. Take a limit $\eta < \gamma$ such that $L_{\eta}$ contains $L_{\bar{\xi}}$ and a witness. In
$L_{\gamma}$, take the canonical well-ordering of $L_{\eta}$ (Section 8; its order certificates exist in
$L_{\gamma}$ by $\Sigma$-recursion), use it to define Skolem functions for $L_{\eta}$, and let $H$ be the
closure of $L_{\bar{\xi}}$ under them. The satisfaction relation of the set $L_{\eta}$ and the enumeration of Skolem
terms are available in the admissible set $L_{\gamma}$, so $H \in L_{\gamma}$, $H$ is countable in $L_{\gamma}$, and $H$
is an elementary substructure of $L_{\eta}$ with $L_{\bar{\xi}} \subseteq H$. Its transitive collapse exists in
$L_{\gamma}$: it is computed by $\Sigma$-recursion on the stage at which each element of $H$ first appears in the
constructible hierarchy, the image of $x$ being the set of images of the elements of $H \cap x$, which appear
earlier. The collapse is a level $L_{\bar{\eta}}$ by (C6), is countable in $L_{\gamma}$, and the collapsing map fixes
$L_{\bar{\xi}}$ pointwise. If $\bar{\eta} \ge \xi_0$, then $L_{\xi_0} \subseteq L_{\bar{\eta}}$ would be countable in
$L_{\gamma}$; so $\bar{\eta} < \xi_0$, and $L_{\bar{\eta}} \models \exists z\, \theta(z, x)$, which gives a witness in
$L_{\xi_0}$.
\end{proof}

\begin{statement}\head{Theorem B} No ordinal below $K$ is \RM{}. Together with Theorem A, the least \RM{} ordinal is
$K$.
\end{statement}

\begin{proof}
Let $\gamma < K$ be a limit above $\omega$. If $\gamma$ is not admissible, it is not \RM{} by Example 3.2(1). Let
$\gamma$ be admissible. Take for $\varphi_0$ the sentence $\forall x\, \exists y\, (y = x)$ with
$q_0 = \emptyset$. Then $L_{\gamma} \models \chi$: $G$ and the clause ``every ordinal is smaller than some limit ordinal'' hold at admissible levels above $\omega$,
and every $L_a$ with $a < \gamma$ is countable in $L_{\gamma}$ by Lemma 7.2. There is no $\Pi_3$-reflecting $\lambda \le \gamma$. By Theorem M, $\gamma$ is
not \RM{}.
\end{proof}

\begin{statement}\head{Lemma 7.3} Let $\gamma > \omega$ be admissible such that every $L_a$ with $a < \gamma$ is countable in $L_{\gamma}$, and
suppose $\gamma$ is \RM{} and not $\Pi_3$-reflecting. Then $\gamma$ is $\Pi_2$-reflecting on $A$.
\end{statement}

\begin{proof}
Otherwise some $\Pi_2$ sentence $\varphi_0(q_0)$ is true in $L_{\gamma}$ and false in $L_{\lambda}$ for every
$\lambda \in A \cap \gamma$ with $q_0 \in L_{\lambda}$. Then $L_{\gamma} \models \chi$, and no $\Pi_3$-reflecting
$\lambda \le \gamma$ with $q_0 \in L_{\lambda}$ satisfies $\varphi_0(q_0)$: the case $\lambda = \gamma$ is excluded because
$\gamma \notin A$. By Theorem M, $\gamma$ is not \RM{}.
\end{proof}

\begin{statement}\head{Lemma 7.4} Every ordinal that is $\Pi_2$-reflecting on $A$ is \RM{}.
\end{statement}

\begin{proof}
Let $\gamma$ be such an ordinal; it is a limit above $\omega$. Let $P$ be a $\Sigma_1$ presentation with
$p \in L_{\gamma}$, $a$ a node in the well-founded part at $\gamma$, and $c$ a predecessor of $a$ at $\gamma$. Then $c$
is in the well-founded part at $\gamma$. Let $\lambda \in A \cap \gamma$ with $p, c \in L_{\lambda}$. By (C5) a descending
sequence from $c$ at $\lambda$ would be one at $\gamma$, so $c$ is in the well-founded part at $\lambda$, and
$\lambda \in C(P, c)$ by Theorem A. Given a $\Pi_2$ sentence true in $L_{\gamma}$, adjoin the constants $p$ and $c$ and
reflect onto $A$: the resulting ordinal lies in $C(P, c)$. So $\gamma$ is $\Pi_2$-reflecting on $C(P, c)$ for every
predecessor $c$ of $a$, and $\gamma \in C(P, a)$.
\end{proof}

\begin{statement}\head{Lemma 7.5} $\gamma_1$ is not $\Pi_3$-reflecting.
\end{statement}

\begin{proof}
$\gamma_1$ is $\Pi_2$-reflecting, hence admissible above $\omega$. Let $S$ be the conjunction of $G$ with
$\forall\varphi\, \forall q\, [\, \neg\Sat_2(\varphi, q) \vee \exists\mu\, (\, q \in L_{\mu} \wedge$ ``$\mu \in A$''
$\wedge L_{\mu} \models \varphi(q)\, )\, ]$, where ``$\mu \in A$'' is expressed positively through a certificate through
$\mu$: given the satisfaction relations of the levels up to $\mu$, the following statement is bounded: $\mu$ is a
limit above $\omega$, and every sentence $\forall u\, \exists v\, \forall w\, \theta$ true in $L_{\mu}$ is true in
$L_b$ for some limit $b$ with $\omega < b < \mu$ whose level $L_b$ contains the constant of the sentence. For such $\mu$ this is equivalent to $\mu \in A$ (reflect
$\psi(q) \wedge G$ for one direction). $S$ is $\Pi_3$, by the computation of Step 4 in the proof of Lemma 5.1:
$\neg\Sat_2$ has the form $\exists j\, \forall k$ and the last disjunct the form $\exists h$, each followed by a
bounded formula, so the second conjunct of $S$ is logically equivalent to a formula with prefix
$\forall \varphi, q\, \exists j, h\, \forall k$ and bounded matrix, and so is its conjunction with $G$. $S$ is true
in $L_{\gamma_1}$, since $\gamma_1$ is $\Pi_2$-reflecting on $A$ and certificates through every $\mu < \gamma_1$ are
elements of $L_{\gamma_1}$, and by (C1) and (C4) any level satisfying it is $\Pi_2$-reflecting on $A$. If $\gamma_1$
were $\Pi_3$-reflecting, a smaller ordinal would satisfy $S$, contradicting the minimality of $\gamma_1$.
\end{proof}

\begin{statement}\head{Theorem C} Let $\gamma > \omega$ be admissible such that every $L_a$ with $a < \gamma$ is countable in $L_{\gamma}$; this
includes every admissible $\gamma$ with $\omega < \gamma \le \gamma_1$. Then $\gamma$ is \RM{} if and only if
$\gamma \in A$ or $\gamma$ is $\Pi_2$-reflecting on $A$. Consequently the least \RM{} ordinal that is not
$\Pi_3$-reflecting is $\gamma_1$.
\end{statement}

\begin{proof}
The equivalence is Theorem A, Lemma 7.3 and Lemma 7.4, with Lemma 7.2 for the range. By Lemmas 7.4 and 7.5,
$\gamma_1$ is \RM{} and not in $A$. An \RM{} ordinal $\gamma < \gamma_1$ that is not in $A$ is admissible, and by Lemma 7.2 every $L_a$ with
$a < \gamma$ is countable in $L_{\gamma}$; so $\gamma$ would be $\Pi_2$-reflecting on $A$ by Lemma 7.3, contradicting
the minimality of $\gamma_1$.
\end{proof}

\begin{statement}\head{Corollary 7.6} Let $\kappa_1$ be the least element of $A$ that is a limit of elements of $A$,
and $\kappa_2$ the least element of $A$ that is $\Pi_2$-reflecting on $A$. Then $K < \kappa_1 < \gamma_1 < \kappa_2$,
and every \RM{} ordinal below $\gamma_1$ is $\Pi_3$-reflecting.
\end{statement}

\begin{proof}
$A \cap \gamma_1$ is unbounded in $\gamma_1$: reflect a trivial sentence with an arbitrary ordinal $\delta < \gamma_1$
as a constant onto $A$. So the $\Pi_2$ sentence ``for every ordinal $\delta$ there is $\nu > \delta$ with
$\nu \in A$'' (certified as in Lemma 7.5, with $G$) is true in $L_{\gamma_1}$. Reflecting it onto $A$ gives an element
of $A$ below $\gamma_1$ that is a limit of elements of $A$, so $\kappa_1 < \gamma_1$. By Lemma 7.1 a
$\Pi_4$-reflecting ordinal is in $A$ and $\Pi_2$-reflecting on $A$, so $\kappa_2$ exists; $\gamma_1 \le \kappa_2$ by
minimality, and $\gamma_1 \ne \kappa_2$ by Lemma 7.5. The last statement is Theorem C.
\end{proof}

\section{Linear orderings}

\textbf{The canonical well-ordering.} The canonical well-ordering of $L$, in which each level is an initial segment,
is $\Sigma_1$-definable, and so $\Delta_1$-definable, over admissible levels (Devlin, Lemma II.3.3). As
with certificates, we fix one form of this definition. An \emph{order certificate through $\xi$} is a certificate $W$
through $\xi$ together with relations $<_\nu$ on $M_\nu$, $\nu \le \xi$, subject to bounded clauses that follow the
construction of the order: $<_0$ is empty, unions are taken at limits, and at a successor the elements of $M_\nu$
come first, in the order $<_\nu$, followed by the new elements, ordered by their least definitions. The relations of
any order certificate are the restrictions to the levels of the canonical well-ordering, and an admissible
$L_{\alpha}$ contains order certificates through every $\xi < \alpha$. So the existence of an order certificate giving
a comparison is $\Sigma_1$, it gives the true comparison when present, and in $L_{\alpha}$ it is present, in one
direction, for any two distinct sets.

\begin{proof}[Self-contained proof]
Augment certificates by the constructible orders: an \emph{order
certificate through $\xi$} is a certificate $W$ through $\xi$ (Section 2) together with a relation $<_\nu$, a set of
pairs from $M_\nu$, for each $\nu \le \xi$, such that $<_0$ is empty, unions are taken at limits, and for $\nu < \xi$
and all $x, y \in M_{\nu+1}$:
$x <_{\nu+1} y$ if and only if either $x, y \in M_\nu$ and $x <_\nu y$, or $x \in M_\nu$ and $y \notin M_\nu$, or
$x, y \notin M_\nu$ and the least defining pair of $x$ precedes that of $y$. Here a pair $(\varphi, s)$ of a
formula code and a tuple $s \in F_\nu$ defines $y$ if it yields $y$ in the sense of (V3), and pairs are ordered by
formula code, a natural number (Section 2), and then by the previous order on tuples, compared by length and then
lexicographically by $<_\nu$. These conditions are bounded on the sets of the certificate,
since ``least defining pair'' quantifies only over formula codes and members of $F_\nu$. By induction on $\nu$, as in
(C1), the relations of any order certificate are the restrictions to the levels of one well-ordering of $L$, of which
each level is an initial segment; this is the canonical well-ordering. As in (C2), an admissible $L_{\alpha}$ contains
order certificates through every $\xi < \alpha$, by $\Sigma$-recursion: at a limit the union of the earlier
relations is taken, and at a successor the next relation is formed by bounded separation from $M_{\nu+1}^2$,
checking the biconditional for every pair. So the existence of an order certificate
giving a comparison is $\Sigma_1$, it gives the true comparison when present, and in $L_{\alpha}$ it is present, in
one direction, for any two distinct sets.
\end{proof}

\begin{statement}\head{Proposition D} Let $\alpha > \omega$ be admissible. Then $\alpha$ is \RM{} if and only if
$\alpha \in C(Q, x)$ for every $\Sigma_1$ presentation $Q$ that defines a linear ordering of $L_{\alpha}$ at $\alpha$
and every node $x$ in the well-founded part of that ordering.
\end{statement}

\begin{proof}
One direction is a restriction of the quantifier. For the other, let $P = (\rho, p)$ be a $\Sigma_1$ presentation,
with $\rho(c, b, p)$ of the form $\exists w\, \theta(c, b, p, w)$, $\theta \in \Delta_0$, and let $a$ be a node in the
well-founded part at $\alpha$. A \emph{history} is a finite sequence
$\langle (c_1, w_1), \dots, (c_n, w_n) \rangle$ with $\theta(c_i, c_{i-1}, p, w_i)$ for each $i$, where $c_0 = a$ by
convention; its \emph{end} is $c_n$, so the empty history has end $a$. With the parameter $(p, a, \omega)$, being a
history is a $\Delta_0$ property, because the edge witnesses are recorded. Define $Q$, with parameter
$(p, a, \omega)$, as follows: histories come before all other sets; between histories, a proper extension comes
before the history it extends, and otherwise histories are compared at the first place where they differ, by the
canonical well-ordering of $L$; sets that are not histories are compared by the canonical well-ordering of $L$.

\emph{The canonical well-ordering is $\Sigma_1$}, through the order certificates described before the proposition.

The formula defining $Q$ is the disjunction of the cases listed above. In each case the conditions on histories are
$\Delta_0$ and a comparison by the canonical well-ordering is $\Sigma_1$; the proper-extension clause for two
histories needs no order certificate. So $Q$ is a $\Sigma_1$ presentation, and at the admissible $\alpha$ it linearly
orders $L_{\alpha}$.

\emph{The empty history is in the well-founded part of $Q$ at $\alpha$.} An infinite descending sequence of histories
in this ordering yields an infinite branch through the tree of histories, hence an infinite descending sequence from
$a$ in the relation defined by $P$ at $\alpha$; there is none. (This is the usual fact that the Kleene--Brouwer
ordering of a well-founded tree over a well-ordered alphabet is a well-ordering.) Every nonempty history lies below the
empty history, and all sets that are not histories lie above it. In fact the whole ordering at $\alpha$ is a
well-ordering, so the proposition also holds with ``well-ordering'' in place of ``linear ordering''.

\emph{Comparison of classes.} We show, by induction on $\mu$ and for all histories $s \in L_{\mu}$:

\begin{quote}
$(\ast)$ $\mu \in C(Q, s)$ implies $\mu \in C(P, \mathrm{end}(s))$.
\end{quote}

Let $\mu \in C(Q, s)$ and let $c$ be a predecessor of $\mathrm{end}(s)$ at $\mu$. Choose a witness $w \in L_{\mu}$ for
this edge and let $t$ be $s$ extended by $(c, w)$. Then $t$ is a history in $L_{\mu}$, and $t$ precedes $s$ in $Q$ at
$\mu$ by the proper-extension clause alone, so this step does not need $\mu$ to be admissible. Since
$\mu \in C(Q, s)$, $\mu$ is $\Pi_2$-reflecting on $C(Q, t)$: every $\Pi_2$ sentence true in $L_{\mu}$ holds in
$L_{\beta}$ for some $\beta < \mu$ in $C(Q, t)$. For such $\beta$, $t \in L_{\beta}$ is still a history with end $c$, so
$\beta \in C(P, c)$ by the induction hypothesis. Thus $\mu$ is $\Pi_2$-reflecting on $C(P, c)$ for every predecessor
$c$ of $\mathrm{end}(s)$, which gives $\mu \in C(P, \mathrm{end}(s))$.

Now if $\alpha$ is in the class of every well-founded node of every $\Sigma_1$ linear ordering, then
$\alpha \in C(Q,$ empty history$)$, and $(\ast)$ gives $\alpha \in C(P, a)$.
\end{proof}

Only the implication $(\ast)$ is proved or needed; a linear ordering has more predecessors than the tree and its
classes can be smaller. The construction does not make the ordering at a smaller level an initial segment of the
ordering at a larger one, and it does not choose one canonical presentation for each order type. A notion with such
coherence requirements would be a different notion, to which this proposition does not apply.

\section{Remarks and questions}

\begin{enumerate}
\item \textbf{The method stops at $\Pi_3$.} In the analogous tree for $\Pi_4$-reflection, the promise left by a
  negative answer is $\Pi_3$, so $\mathrm{corr}(s)$ is $\Pi_3$ and $\tau$ is no longer $\Pi_2$; the classes of
  Definition 3.1 only reflect $\Pi_2$ sentences. The argument must fail there: $K$ is \RM{} but is not
  $\Pi_2$-reflecting on the $\Pi_4$-reflecting ordinals.
\item \textbf{A related result.} Richter and Aczel show that the supremum of the closure ordinals of effective
  $\Pi^0_2$ operators on subsets of $\omega$ (non-monotone operators allowed, with accumulating iteration) is $K$.
  We expect a second proof of the lower bound in Theorem B from a tree that guesses the closed stage of a universal
  $\Pi^0_2$ inductive definition; this paper does not give that proof. The proof given here uses Theorem M.
\item \textbf{Question.} Is \RM{} equivalent to a recursive form of the canonical-function or normal-ideal definition
  of greatly Mahlo cardinals, with suitable coherence requirements on the presentations? (See the end of Section 8.)
  For cardinals, Arai relates his Mahlo classes of degree $\sigma^+$ to the greatly Mahlo cardinals of Baumgartner,
  Taylor and Wagon (Proof theory of weak compactness, Proposition 2.10); that would be the place to start.
\item \textbf{Question.} Theorem C assumes that every $L_a$ with $a < \gamma$ is countable in $L_{\gamma}$, which fails at
  sufficiently large admissible $\gamma$. What is \RM{} there?
\item \textbf{Question.} One level down, with ``admissible limit of'' in place of ``$\Pi_2$-reflecting on'' and
  recursively Mahlo in place of $\Pi_3$-reflecting, the analogous statements appear to hold by the analogous argument.
  We have only sketched this. The provision ``limit of limit ordinals'' used for making room in Lemma 6.2 has no
  analogue two levels down.
\item \textbf{Question.} By the remark after Definition 3.3, \RM{} only concerns degrees below $\alpha^+$. Nothing in
  this paper addresses degrees from $\alpha^+$ on. They would need presentations defined over levels above
  $L_{\alpha}$, and we do not know whether a natural notion of that kind matches a stronger reflection property or a
  stability property.
\end{enumerate}

\section{Verification status}

\begin{center}\begin{tabular}{p{0.24\linewidth}p{0.70\linewidth}}
\toprule
Part & Status \\
\midrule
Lemma 4.1 & The abstract descent theorem is formalised in Lean 4 with Mathlib, and the axiom audit reports no axioms
for it. Its applications to $L$, the tree, Theorems A--C and Proposition D are written mathematics. \\
Theorems A--C, Proposition D & Written proofs, not checked by a person. Theorem A is a version of a known
argument. \\
Facts (C1)--(C6) & (C2), (C5), (C6) are standard; Section 2 gives sources. (C1), (C3), (C4) are stated in the form used, with their
justifications written out in Section 2; they are not taken from a source. \\
References & The Richter--Aczel statements were checked against the 1973 Oslo preprint, and the Arai statements
against the arXiv versions cited. Gostanian's theorem is cited from his paper (Section 2, application (viii); Corollary 3.2), for countable admissible $\alpha$ with every element of
$L_{\alpha}$ countable in $L_{\alpha}$ and for well-orderings of subsets of $\alpha$; Section 3.3 transfers it to
well-orderings with field contained in $L_{\alpha}$. The papers of Rathjen and of Baumgartner, Taylor and Wagon were not read in full. \\
Novelty of B, C, D & Theorem A is Arai's argument. A search of the sources listed found no statement of Theorems B
and C or of Proposition D; Arai's Theorem 1.3 is a proof-theoretic counterpart of B. Rathjen's paper was not read.
Not established. \\
\bottomrule
\end{tabular}\end{center}

No human specialist has reviewed the proofs. The passages a specialist is most likely to question are the soundness of
certificates (C1), the step in Lemma 7.1 that reflects a $\Pi_3$ sentence to a limit level and reads it there with each block of like
quantifiers merged into one, and the role of the
canonical well-ordering of $L$ in Section 8: it is needed only to compare sets at the admissible $\alpha$, while the
induction at smaller ordinals uses only the clause that puts a proper extension of a history before that history,
which needs no order certificate. Beyond correctness, the open matter is whether Definition 3.3 is the notion one wants (Sections 1, 8
and Question 9.3).

\section*{Appendix: the Lean statement of Lemma 4.1}

The ordinals are abstracted to a type \texttt{Pt} with a well-founded order, nodes to a type \texttt{N}, and tests to a
type \texttt{T}.
\texttt{holds t μ} means the test \texttt{t} is true at \texttt{μ} with its constants present; \texttt{pred μ c a}
means \texttt{c} is a predecessor of \texttt{a} at \texttt{μ}.

\begin{verbatim}
def Refl (holds : T → Pt → Prop) (μ : Pt) (X : Pt → Prop) : Prop :=
  ∀ t, holds t μ → ∃ β, β < μ ∧ X β ∧ holds t β

-- the classes of Definition 3.1, by well-founded recursion on μ
theorem cls_iff (μ : Pt) (a : N) :
    cls holds pred ok μ a ↔
      ok μ a ∧ ∀ c, pred μ c a → Refl holds μ (fun β => cls holds pred ok β c)

-- Lemma 4.1
theorem not_cls_of_family
    (hZ : ∀ δ c, Z δ c → ∃ d t, pred δ d c ∧ holds t δ ∧ ∀ β, β < δ → holds t β → Z β d)
    {δ : Pt} {c : N} (hc : Z δ c) : ¬ cls holds pred ok δ c
\end{verbatim}

Here \texttt{ok μ a} stands for ``$\mu$ is a limit above $\omega$ with $p, a \in L_{\mu}$''. To apply the theorem, a
test is a $\Pi_2$ sentence together with the node $d$, and \texttt{holds t β} means that $\beta$ is a limit above
$\omega$, $d$ and the constants of the sentence are in $L_{\beta}$, and the sentence is true in $L_{\beta}$. Only the descent is formalised; the
construction of the family $Z$ in Section 6 is not.

\section*{References}

\begin{itemize}
\item J. P. Aguilera, The order of reflection, arXiv:1906.11769 (Section 2).
\item T. Arai, Iterating the recursively Mahlo operations, arXiv:1005.1987 (Definition 1.1, Lemma 1.2, Theorem 1.3).
\item T. Arai, Lectures on Ordinal Analysis, arXiv:2304.00246v1, \S3.2 (Lemma 3.33, Remark 3.34).
\item T. Arai, Proof theory of weak compactness, arXiv:1111.0462 (Proposition 2.10).
\item J. Barwise, \emph{Admissible Sets and Structures}, Springer, 1975.
\item J. E. Baumgartner, A. D. Taylor and S. Wagon, On splitting stationary subsets of large cardinals,
  \emph{Journal of Symbolic Logic} 42 (1977), 203--214.
\item K. J. Devlin, \emph{Constructibility}, Springer, 1984.
\item R. Gostanian, The next admissible ordinal, \emph{Annals of Mathematical Logic} 17 (1979), 171--203.
\item M. Rathjen, Proof theory of reflection, \emph{Annals of Pure and Applied Logic} 68 (1994), 181--224.
\item W. Richter and P. Aczel, Inductive definitions and reflecting properties of admissible ordinals, in
  \emph{Generalized Recursion Theory}, North-Holland, 1974, 301--381. (Checked in the Oslo preprint, 1973: Theorem 1.8,
  Theorem A, Corollary 9.3.)
\end{itemize}

\end{document}
